\originalchapter{多维波与界面折射}\label{ch:interface}
\originalsection{波矢与二维正常模}
鼓膜上的不同节点线围出不同振动区域, 说明二维波还需要记录传播方向. 在平面上, 波矢 $\boldsymbol k=(k_x,k_y)$ 给相位 $\boldsymbol k\cdot\boldsymbol r-\omega t$. 等相位线垂直于 $\boldsymbol k$, 是平面波的波前. 无限二维珠网中, 横、纵间距 $a_x,a_y$, 每个质量 $m$ 仅沿 $z$ 运动, 连结线张力为 $T_x,T_y$. 最近邻的小斜率回复作用给
\[
\omega^2=\frac{4T_x}{ma_x}\sin^2\frac{k_xa_x}2+
\frac{4T_y}{ma_y}\sin^2\frac{k_ya_y}2.
\]
固定 $\omega$ 时可有许多方向与波数组合. 在有限 $N_x\times N_y$ 珠网四边固定, $k_x=n_x\pi/[(N_x+1)a_x]$, $k_y=n_y\pi/[(N_y+1)a_y]$, 空间模为两正弦的乘积, $1\le n_x\le N_x$, $1\le n_y\le N_y$. 边界离散化波数, 但不同方向仍可能给相同频率, 称为简并.

等间距、等张力网的长波极限为膜方程 $\psi_{tt}=c_m^2(\psi_{xx}+\psi_{yy})$, $c_m^2=\mathcal T/\rho_s$, 其中面张力 $\mathcal T=T/a$ 单位 $\mathrm{N/m}$, 面密度 $\rho_s=m/a^2$ 单位 $\mathrm{kg/m^2}$. 长方形膜边长 $L_x,L_y$ 四边固定, 模形状为 $\sin(n_x\pi x/L_x)\sin(n_y\pi y/L_y)$, 频率平方为 $c_m^2[(n_x\pi/L_x)^2+(n_y\pi/L_y)^2]$. 在正方形上交换 $n_x,n_y$ 不改频率, 正是几何对称导致的简并.

三维标量波满足 $\psi_{tt}=c^2\nabla^2\psi$. 对刚性封闭长方体声腔, 正确的是压力法向导数为零, 所以声压模为 $\cos(n_x\pi x/L_x)\cos(n_y\pi y/L_y)\cos(n_z\pi z/L_z)$, 各 $n_j\ge0$ 且不全为零. 若是标量场本身在壁面为零的 Dirichlet 模型, 才改用三个正弦且各 $n_j\ge1$. 线性声学动量式 $\rho_0\boldsymbol v_t=-\nabla p'$ 与刚壁 $v_n=0$ 给非零频振动的 $\partial_np'=0$. 声腔频率为 $\omega=c_s\pi\sqrt{(n_x/L_x)^2+(n_y/L_y)^2+(n_z/L_z)^2}$. 两者频率结构相似, 物理边界不同. 空间无界球面向外波为 $\psi=F(r-ct)/r$, 在 $r>0$ 代入径向拉普拉斯算符即可验证; 振幅 $1/r$ 对应球面面积增长.

\originalsection{相位匹配与斯涅尔定律}
平直静止界面上, 入射、反射、透射的场必须在每个界面点、每一时刻满足边界条件. 各波若具有不同时间频率或不同切向波数, 指数因子不能在整个界面相互匹配. 故 $\omega_i=\omega_r=\omega_t$, $\boldsymbol k_{\parallel i}=\boldsymbol k_{\parallel r}=\boldsymbol k_{\parallel t}$. 同介质中法向波数反号, 所以反射角等于入射角. 各向同性介质内 $|\boldsymbol k|=n\omega/c$, 因而
\begin{equation}\label{eq:snell}
n_1\sin\theta_i=n_2\sin\theta_t.
\end{equation}
角度从界面法线量起. 频率不变而波长在介质中变为真空波长的 $1/n$. 相位匹配决定方向, 力或场的边界条件决定振幅; 两步不能混为一谈.

当 $n_1>n_2$, 入射角超过 $\theta_c=\arcsin(n_2/n_1)$, 透射法向波数满足 $k_{zt}^2=(n_2\omega/c)^2-k_\parallel^2<0$. 取向介质深处衰减的分支 $k_{zt}=\ii\kappa$ 后, 仍有倏逝场 $\e^{-\kappa z}$, 而非“介质内完全没有场”. 对无损半无限介质, 透入深处的平均法向能流为零, 反射能流等于入射能流. 全反射的相位变化影响偏振, 是波动光学能从几何光学之外补充的信息.

\originalsection{菲涅耳振幅与能流}
设界面为 $z=0$, 入射面为 $xz$ 平面, 两介质无自由表面电流、电荷. 边界要求切向 $\boldsymbol E,\boldsymbol H$ 连续, 法向 $\boldsymbol D,\boldsymbol B$ 连续. 记 $Z_j=\sqrt{\mu_j/\epsilon_j}$ 为电磁阻抗. s偏振电场沿 $y$; p偏振采用入射单位方向 $(\cos\theta_i,0,-\sin\theta_i)$, 反射方向 $(\cos\theta_i,0,\sin\theta_i)$, 透射方向 $(\cos\theta_t,0,-\sin\theta_t)$. 此约定使正入射反射电场幅度的符号与固定 $x$ 方向比较一致.

s偏振的连续条件为 $1+r_s=t_s$, $(\cos\theta_i/Z_1)(1-r_s)=(\cos\theta_t/Z_2)t_s$. p偏振为 $\cos\theta_i(1+r_p)=\cos\theta_t t_p$, $(1-r_p)/Z_1=t_p/Z_2$. 解这两个二元系统得到
\begin{align*}
r_s&=\frac{Z_2\cos\theta_i-Z_1\cos\theta_t}{Z_2\cos\theta_i+Z_1\cos\theta_t},&
t_s&=\frac{2Z_2\cos\theta_i}{Z_2\cos\theta_i+Z_1\cos\theta_t},\\
r_p&=\frac{Z_2\cos\theta_t-Z_1\cos\theta_i}{Z_2\cos\theta_t+Z_1\cos\theta_i},&
t_p&=\frac{2Z_2\cos\theta_i}{Z_2\cos\theta_t+Z_1\cos\theta_i}.
\end{align*}
这些是电场振幅系数. 传播透射支上功率系数为 $R=|r|^2$, $\mathcal T=(Z_1/Z_2)(\cos\theta_t/\cos\theta_i)|t|^2$, 并且 $R+\mathcal T=1$. 对非磁性介质 $Z\propto1/n$. 正入射时 $r=(n_1-n_2)/(n_1+n_2)$, $t=2n_1/(n_1+n_2)$. 从较小折射率入射较大折射率时 $r<0$, 反射电场发生 $\pi$ 相移.

非磁性无损介质的 p反射在布儒斯特角消失. 令分子为零, 得 $\cos\theta_t/\cos\theta_i=n_2/n_1$, 配合式 \eqref{eq:snell} 得 $\theta_i+\theta_t=\pi/2$, $\tan\theta_B=n_2/n_1$. 两介质相同的退化情况本来就没有反射. 在布儒斯特角 $\mathcal T=1$ 是功率陈述, 电场幅度 $t_p=\cos\theta_i/\cos\theta_t=n_1/n_2$ 一般不是1. s偏振通常仍反射, 所以用此角可从非偏振光的反射部分获得线偏振光.
